% Generated draft: requires mathematical and source review.
\documentclass[11pt]{article}
\usepackage[margin=25mm]{geometry}
\usepackage{fontspec}
\setmainfont{Latin Modern Roman}
\usepackage{amsmath,amssymb,mathtools,unicode-math}
\setmathfont{Latin Modern Math}
\usepackage{enumitem,needspace,xurl,hyperref}
\setlength{\parindent}{0pt}
\setlength{\parskip}{5pt}
\setlength{\emergencystretch}{3em}
\newcommand{\N}{\mathbb{N}}
\newcommand{\Z}{\mathbb{Z}}
\newcommand{\Q}{\mathbb{Q}}
\newcommand{\R}{\mathbb{R}}
\newcommand{\C}{\mathbb{C}}
\newcommand{\F}{\mathbb{F}}
\newcommand{\A}{\mathbb{A}}
\newcommand{\PP}{\mathbb{P}}
\newcommand{\E}{\mathbb{E}}
\newcommand{\eps}{\varepsilon}
\newcommand{\dd}{\,\mathrm{d}}
\begin{document}
\section*{UnsolvedMath ID: 2308014}
\textbf{Source problem code:} catalogue entry
\subsection{Definitions and mathematical statement}
Let $\mathbb T=\{e^{it}:0\leq t<2\pi\}$ with normalized arc measure, and let $\varphi\in L^\infty(\mathbb T)$ be a symbol, defined up to equality almost everywhere. F. Holland asks which such symbols generate positive Hankel operators.

The source does not define the Hankel convention or the meaning of positive. In one common \emph{linear matrix} convention, write
\[
\widehat\varphi(n)=\frac1{2\pi}\int_0^{2\pi}\varphi(e^{it})e^{-int}\,dt,
\qquad
(\Gamma_\varphi)_{jk}=\widehat\varphi(j+k+1),\quad j,k\geq0,
\]
with the corresponding bounded Hankel operator acting on $\ell^2(\Z_{\geq0})$. Positivity then means positive semidefiniteness:
\[
\sum_{j,k\geq0}\widehat\varphi(j+k+1)c_k\overline{c_j}\geq0
\quad\text{for every finitely supported complex sequence }(c_j),
\]
where the quadratic form is real. The resulting characterization problem is to describe the $L^\infty$ symbols satisfying this condition in more informative function-theoretic terms, taking into account that symbols with the same relevant Fourier coefficients generate the same operator.

This convention is conditional, not a recovered source definition. Other Hardy-space conventions use negative Fourier indices or a different shift. Moreover, a Hankel map written from $H^2$ to its orthogonal complement needs an identification of its target with its source before operator positivity is meaningful, and anti-linear Hankel operators require a different convention. The record provides none of these choices. The supplied summary calls the problem open, notes the intricacy of positive-Hankel-generating symbols, and reports no definitive single resolution located; it supplies no precise characterization theorem.
\subsection{Short English statement}
Holland asks for a characterization of bounded symbols producing positive Hankel operators. A standard matrix version is made explicit, but the source omits the operator convention and target-space identification needed to give positivity a unique meaning.
\subsection{Sources}
\begin{itemize}
\item Local UnsolvedMath registry, permanent ID 2308014, exact source code catalogue entry, ``Research Problems in Function Theory --- Problem 8.14''; status open, classification OPEN-TRIAGE.
\item Hayman, \emph{Research Problems in Function Theory} (2018), Problem 8.14, attributed to F. Holland, \url{https://arxiv.org/abs/1809.07200}. The record identifies extraction from public arXiv source TeX, accessed July 29, 2026; the 2018 update was retained for status triage.
\item The supplied research summary labels the characterization open and says no definitive single resolution was located. No independent authoritative citation, operator definition, or precise partial theorem is supplied.
\end{itemize}
\end{document}
